% !TEX root = ../algebra_02.tex

\section{Строение конечных и конечно порождённых абелевых групп (формулировки)}

\subsection{Конечные абелевы группы}

\begin{Theorem}{Основная теорема о конечных абелевых группах}{}
	Пусть $G$ --- конечная абелева группа.
	Тогда $G$ изоморфна прямому произведению циклических групп примарных порядков:
	\[
	G \cong \prod \mathbb{Z}/p_i^{n_i}\mathbb{Z},
	\]
	где $p_i$ --- простые числа, $n_i \in \mathbb{N}$.
	Такое представление единственно с точностью до перестановки сомножителей.
\end{Theorem}

\begin{Example}{Возможные абелевы группы порядка 200}{}
	Пусть $|G| = 200 = 2^3 \cdot 5^2$.
	Чтобы найти все возможные (с точностью до изоморфизма) абелевы группы такого порядка, переберем разбиения степеней простых множителей.

	Варианты для $p=2$ (степень 3):
	\begin{itemize}
		\item $2, 2, 2 \implies (\mathbb{Z}/2\mathbb{Z})^3$ 
		\item $2^2, 2 \implies (\mathbb{Z}/4\mathbb{Z}) \times (\mathbb{Z}/2\mathbb{Z})$ 
		\item $2^3 \implies \mathbb{Z}/8\mathbb{Z}$ 
	\end{itemize}

	Варианты для $p=5$ (степень 2):
	\begin{itemize}
		\item $5^2 \implies \mathbb{Z}/25\mathbb{Z}$ 
		\item $5, 5 \implies (\mathbb{Z}/5\mathbb{Z})^2$ 
	\end{itemize}

	Комбинируя их, получаем список всех возможных групп.
	Например, при множителе $\mathbb{Z}/25\mathbb{Z}$:
	\begin{align*}
		1) \quad & (\mathbb{Z}/2\mathbb{Z})^3 \times (\mathbb{Z}/25\mathbb{Z}) \\
		2) \quad & (\mathbb{Z}/4\mathbb{Z}) \times (\mathbb{Z}/2\mathbb{Z}) \times (\mathbb{Z}/25\mathbb{Z}) \\
		3) \quad & (\mathbb{Z}/8\mathbb{Z}) \times (\mathbb{Z}/25\mathbb{Z}) \cong \mathbb{Z}/200\mathbb{Z}
	\end{align*}
\end{Example}

\subsection{Конечнопорожденные абелевы группы}

\begin{Definition}{Конечнопорожденная группа}{}
	Группа $G$ называется конечнопорожденной, если существует конечный набор элементов $g_1, \dots, g_s \in G$ такой, что они порождают всю группу: $G = \langle g_1, \dots, g_s \rangle$.
\end{Definition}

\begin{Definition}{Подгруппа кручения (периодическая часть)}{}
	Пусть $G$ --- абелева группа. Множество всех её элементов конечного порядка образует подгруппу $T(G)$, которая называется подгруппой кручения:
	\[
	T(G) = \{ g \in G \mid \operatorname{ord}(g) < \infty \} < G \text{ }.
	\]
\end{Definition}

\begin{Theorem}{Структура конечнопорожденной абелевой группы}{}
	Пусть $G$ --- конечнопорожденная абелева группа.
	Тогда в $G$ существует свободная подгруппа $H$ такая, что $G$ является внутренним прямым произведением $H$ и $T(G)$:
	\[
	G = H \times T(G),
	\]
	причем:
	\begin{enumerate}
		\item $T(G)$ --- конечная группа;
		\item $H \cong \mathbb{Z}^r$, где $r \ge 0$.
	\end{enumerate}
	Число $r$ называется рангом группы $G$ и является её инвариантом (не зависит от выбора $H$).
\end{Theorem}

\begin{Example}{Неединственность выбора свободной части $H$}{}
	Рассмотрим группу $G = \mathbb{Z} \times (\mathbb{Z}/2\mathbb{Z})$.
	Её подгруппа кручения определяется однозначно:
	\[
	T(G) = \{ (0, \bar{0}), (0, \bar{1}) \} \text{ }.
	\]
	В качестве свободной подгруппы $H \cong \mathbb{Z}$ можно взять:
	\[
	H = \{ (a, \bar{0}) \mid a \in \mathbb{Z} \} = \langle (1, \bar{0}) \rangle \text{ }.
	\]
	Однако этот выбор не уникален. Мы можем взять подгруппу $H' = \langle (1, \bar{1}) \rangle$, которая также изоморфна $\mathbb{Z}$ и удовлетворяет условию $G = H' \times T(G)$.
	Это показывает, что $H$ в разложении определяется неоднозначно.
\end{Example}
